\section{The People Entrusted to Our Care}

A responsibility for persons brings its own duties before judgment. In Luke's account of the vigilant servants, Peter asks whether the parable is addressed especially to the disciples or to everyone. Christ answers with the faithful steward appointed over a household to give its members their portion in due season (Lk 12:41--44). The steward receives authority to provide for the household. He must supply what its members need at the time they need it. His account to the master will concern the fulfillment of that duty.

The unfaithful steward reverses this order. Deciding that the master delays, he beats the servants and abandons himself to appetite (Lk 12:45--46). Those whose care justified his authority become objects of his convenience or anger. He treats the household as though the master's absence had made it his possession. The sudden return reveals the falsehood. An entrusted office cannot be transformed into a private right to use others, however long the misuse has continued or however little resistance its victims have been able to offer.

The Lord's warning reaches familiar forms of responsibility. A parent, teacher, employer, superior, or public official makes decisions that affect people who may have less freedom to resist. The responsibility may be modest in the world's eyes and still be decisive for someone under it. We should therefore examine the place from which our choices reach another person's daily life. The relevant question is the good entrusted to us: whether those within our care received justice, truthful guidance, patient attention, and the protection which our position made possible.

Authority can hide faults from the person exercising it. Those who depend upon us may remain silent because correcting us feels dangerous or useless. Their compliance can then be mistaken for agreement, and our freedom from criticism for proof of wisdom. The final judgment will not accept that mistake. Christ knows whether peace came from just conduct or from another person's fear of speaking. Remembering his judgment should make a person in authority more willing to hear an unwelcome account, seek counsel, and permit the correction which pride would prefer to exclude.

Admitting a fault before someone entrusted to us can seem to diminish the authority required for our work. Defending an error, however, leaves the person in our care bearing its consequences. A teacher who discovers that an explanation was false owes the hearers a correction. A parent who acted unjustly can acknowledge the act and repair it while continuing to guide the child. The duty to govern includes accepting responsibility for errors made in its exercise.

James makes the account of teachers especially serious. He warns against becoming teachers in great numbers because they receive a stricter judgment, and immediately speaks of the failings of the tongue (Jas 3:1--2). Religious instruction brings this warning close to anyone who speaks publicly about God. The capacity to explain doctrine is a gift for the hearer's good. It calls for accuracy, preparation, and submission to the truth taught. A teacher should fear misleading another person through negligence, making conjecture sound certain, or using sacred teaching to magnify himself.

The warning also requires a life willing to receive the command it explains. We can teach patience while excusing our own cruelty, speak of poverty while indulging greed, or exhort repentance while refusing correction. The discrepancy does more than expose personal weakness. It can make the hearer distrust the truth because its messenger appears to use it only against others. Paul's examination of the religious teacher in Romans returns the command to the one who teaches it: the person who instructs another must also receive instruction himself (Rom 2:17--24).

Those entrusted with teaching still have a duty to instruct while seeking their own amendment. We can study what we do not understand, distinguish the Church's doctrine from our own opinion, and correct an error we have spread. The warning about stricter judgment also gives a reason to examine our conduct before teaching a commandment. Where it exposes a failure, we should seek forgiveness and amendment rather than hide the teaching from those who need it.

Other forms of authority concern goods and labor. James speaks with dreadful directness about wages withheld from those who reaped the fields: the workers' cry has reached the Lord (Jas 5:1--6). Here judgment enters the apparent distance between prosperity and the lives that helped produce it. A person may enjoy abundance while concealing from himself what was denied to someone with less power. The apostle refuses that concealment. God hears the claim which the prosperous person has found ways not to hear.

This passage asks us to bring our practical arrangements before conscience. An employer's religious devotion cannot make dishonest withholding just. A person administering another's goods cannot treat those goods as available for private advantage. A household cannot purchase its comfort through a deliberate refusal of obligations owed to someone dependent upon it. Particular questions can require competent advice, but the need for advice must not become a permanent shelter from an obligation already clear. Christ's judgment reaches the decision by which we preferred our advantage to another person's due.

Works of generosity do not erase such duties. Giving alms can be a true good, yet a gift offered publicly does not authorize the retention of what belongs to an unpaid worker. The person praised as a benefactor may still need to restore what he wrongfully kept. Annual meditation can uncover this divided account, in which selected acts of generosity have made us unwilling to examine the way our resources are obtained or administered. Repentance makes the whole relation to goods answerable to justice and charity, including the parts from which we derive little admiration.

An entrusted duty can also be neglected through the desire to avoid discomfort. A person responsible for guidance may decline to correct a dangerous habit because correction risks resentment. Patience is needed in deciding when and how to speak. If the delay leaves someone exposed to an avoidable danger, however, the reason for delay needs examination. We may have preferred freedom from conflict to the protection which our responsibility required us to provide.

The opposite temptation is to make ourselves responsible for what belongs to another person's freedom. Parents can teach and give example; they cannot compel an adult child's interior assent. Teachers can explain; they cannot guarantee that knowledge will become obedience. Chrysostom's treatment of the talents helps here: the servant must offer the counsel entrusted to him, while the result is left to the master (LXXVIII.2). The distinction guards authority from both negligence and domination. We must do our part faithfully without claiming the power to possess another person's response.

Some people carry deep grief over failures within those entrusted to them. Remembering judgment may awaken a truthful need for repentance, but it can also tempt them to assume every later evil as their own deliberate act. Christ will distinguish the care given, the omission freely chosen, the limitation honestly encountered, and the independent choice of another person. We should acknowledge an actual failure and repair it where possible. Where grief has made this distinction difficult, sound counsel can help us identify what was within our responsibility and entrust the rest to God.

James places patient endurance beside his warning to the unjust. The wronged are to strengthen their hearts for the Lord's coming; they are also warned against resentful complaint against one another, because the Judge is near (Jas 5:7--9). Those who have suffered through another person's misuse of power remain called to charity. They may seek protection and a just remedy without allowing injury to become permission for cruelty toward someone else. Awaiting Christ means entrusting the full account to him while remaining obedient in the part of life still placed in our own hands.

Stewardship may require perseverance through a long period of little apparent progress. A teacher can continue preparing carefully and attending to a pupil who learns slowly. Someone administering a difficult trust can continue seeking counsel and giving a fair account. Such duties need examination as well as encouragement: persistence in a method does not by itself prove that the people concerned are receiving good care. The faithful steward attends to their needs while remaining answerable to the master.

